\begin{theorem}[DP38: the first primitive Poincar\'e coefficients]\label{DP38}
For every $p$ and $h\ge2$, over $W(\F_q)$ there are integral Teichm\"uller coordinates $S,Y_1,\ldots,Y_d$ on the Honda line and its actual dual, $d=h-1$, and a normalized integral Poincar\'e frame in which every rational compatible splitting satisfies
\[
 [S^pY_1]B=\frac1p,\qquad [SY_d^p]B=-\frac1p.
\]
These coefficients are unchanged by a rational bilinear gauge and change only by integral amounts under an axes-normalized integral unit frame. The frame is invariant for the true action $(S,Y_i)\mapsto(\zeta S,\zeta^{-p^i}Y_i)$. Its unique invariant logarithmic splitting has no positive term below degree $p+1$, and
\begin{equation}\label{eq:primitive-first}
 (\log B)_{p+1}=B_{p+1}=\frac{S^pY_1-SY_d^p}{p}.
\end{equation}
The same first homogeneous term holds in the normalized logarithm coordinates. Consequently
\[
 D_G^pD_1 B(0)=(p-1)!,\qquad D_GD_d^p B(0)=-(p-1)!.
\]
Both values are units, although both operators are divisible by $p$ on integral formal functions. In the group coordinates,
\[
 w_t(B)=-\infty\quad\left(0<t<\frac1{p+1},\ t\in\Q\right),\qquad
 \liminf_{|\alpha|\to\infty}\frac{v_p(B_\alpha)}{|\alpha|}\le-\frac1{p+1}.
\]
For this invariant splitting, odd $p$, $h>2$, and any integral $U$ with unit constant term, the last-dual logarithmic path $x=t$, $y_d=\epsilon t$, $y_i=0$ for $i<d$, gives
\begin{equation}\label{eq:primitive-path-unit}
 v_p\left((p+1)![t^{p+1}]\,
 U(E_G(t),E_H(\epsilon te_d))
 \exp\left(\tfrac12\log B(E_G(t),E_H(\epsilon te_d))\right)\right)=0
\end{equation}
for every unit $\epsilon$. This is an exact finite-order assertion, not an asymptotic upper bound on every selected theta jet or a shorter-chain impossibility theorem.
\end{theorem}
\begin{proof}
We give an explicit display calculation, working compatibly over $\mathcal O/p^n$, where $\mathcal O=W(\F_q)$, and then passing to the inverse limit. Let $P$ be free on $e_0,\ldots,e_d$ over the Witt ring, with normal decomposition
\[
 T=We_0,\qquad L=\bigoplus_{j=1}^dWe_j,\qquad Q=L\oplus IT.
\]
With $e_h=e_0$, specify the Frobenius-semilinear operators by
\[
 Fe_0=e_1,\qquad V^{-1}e_j=e_{j+1},\qquad Fe_j=pe_{j+1}\quad(1\le j\le d).
\]
On $IT$ the relation is $V^{-1}(V\xi\,e_0)=\xi Fe_0$. The linearization of $V^{-1}\oplus F$ on $L\oplus T$ is a cyclic permutation and hence an isomorphism. The actual dual display has
\[
 L'=We_0^*,\quad T'=\bigoplus_{j=1}^dWe_j^*,\quad
 V'^{-1}e_0^*=e_1^*,\quad F'e_0^*=pe_1^*,\quad
 F'e_j^*=e_{j+1}^*\ (j\ge1).
\]
These formulas satisfy the pairing identities of Zink, Definition~19 and (22), p.~143. The $h$-fold Frobenius cycles send a basis vector to $p^{h-1}$ times itself in the first display and to $p$ times itself in the dual, with the corresponding Frobenius on coefficients. Both nilpotence conditions hold for $h\ge2$.

By Zink's Theorem~81 and (145)--(146), p.~204, the formal groups are the cokernels of $V^{-1}-\id$. The maps represented by
\[
 \sigma(S)=[S]e_0,\qquad \tau(\mathbf Y)=\sum_{j=1}^d[Y_j]e_j^*
\]
give integral formal coordinates: their tangent maps to $P/Q$ and $P'/Q'$ are identities, so the formal inverse function theorem, proved coefficientwise by its unit linear term, applies. On the special fibre $p[S]=V[S^p]$, and the display relations carry this around the cycle to $[S^q]e_0$. Thus $[p]_G(S)\bmod p=S^q$. The linear coefficient is $p$, so Hazewinkel's integral Lubin--Tate uniqueness theorem, \S5, p.~77 identifies this line with $G_h$. The second display is its actual Cartier dual. No equality of its entire coordinate logarithm with DP32's chosen formula is assumed in the calculation that follows.

Use the evaluation pairing of the displays. The scalar-one section represented by $(1,\sigma(S),\tau(\mathbf Y))$ in Zink's Proposition~106, p.~226 gives a normalized integral frame. If $\kappa_G,\kappa_H$ are the defects of these lifts, defined by
\[
 (V^{-1}-\id)\kappa_G=\sigma(S)+\sigma(S')-\sigma(F_G(S,S'))
\]
and its dual analogue, its cocycles from (218), p.~239 are
\begin{equation}\label{eq:display-cocycles}
 C_G=\operatorname{hex}\langle V^{-1}\kappa_G,\tau(\mathbf Y)\rangle,
 \qquad C_H=\operatorname{hex}\langle\sigma(S),\kappa_H\rangle^{-1}.
\end{equation}
Here $\operatorname{hex}$ is the Artin--Hasse map (210)--(211), p.~236. The inverse is the minus sign in Zink's additively written multiplicative group. Indeed, the quotient relation in Proposition~106 replaces a lift differing by a defect with its scalar multiplied by the corresponding pairing. Thus these signs agree with $C_G=B(F_G)/(BB)$, and similarly for $C_H$. DP35 supplies rational compatible splittings.

Put
\[
 C_p(S,S')=\frac{(S+S')^p-S^p-S'^p}{p},\qquad
 [S]+[S']-[S+S']=V\xi.
\]
The first two ghost polynomials give $\xi_0=-C_p(S,S')$. The integral Witt addition polynomials are weighted homogeneous, so $\xi_i$ has ordinary degree $p^{i+1}$, and every component of $F^r\xi$ has degree at least $p^{r+1}$. This assertion uses this homogeneous carry vector; it is not an assertion about Frobenius and arbitrary coefficient-degree filtrations. The explicit element
\[
 \kappa_0=-V\xi\,e_0-\sum_{j=1}^dF^{j-1}\xi\,e_j
\]
satisfies the exact telescoping identity
\[
 (V^{-1}-\id)\kappa_0=V\xi\,e_0-F^d\xi\,e_0.
\]
The last term has degree at least $p^{d+1}=q$. Modulo the degree-$q$ ideal the group law is therefore ordinary addition. Moreover the true defect equals $\kappa_0$ to that precision: reduce Zink's exact sequence (146) to this nilpotent test algebra and use its injectivity. No inverse of an uncomputed operator is assumed to preserve degree.

The $e_1$ coefficient of $V^{-1}\kappa_0$ is the Witt vector $w=-\xi$. Multiplication by $[Y_1]$ sends its $i$th Witt component to $w_iY_1^{p^i}$; no componentwise-negation formula is assumed. In the Artin--Hasse product only the zeroth component contributes linearly in $Y_1$, with coefficient one, and $w_0=-\xi_0$. Thus
\[
 [Y_1]\log C_G=C_p(S,S')+O((S,S')^q).
\]
Writing $b=\log B$, extraction of $S^{p-1}S'Y_1$ from $\delta_G b$ uses ordinary addition and gives $p[S^pY_1]b=1$. Since $b$ has positive degree on both sides, exponentiation cannot alter its $Y$-linear coefficient. This proves $[S^pY_1]B=1/p$.

For the second coefficient restrict both dual inputs to the last coordinate. Again write $[Y]+[Y']-[Y+Y']=V\xi$. This time
\[
 \kappa'_0=-V\xi\,e_d^*-\xi\,e_0^*,\qquad
 (V'^{-1}-\id)\kappa'_0=V\xi\,e_d^*-F\xi\,e_1^*.
\]
The endpoint has degree at least $p^2$. Addition on this source axis is ordinary to that precision, and the true defect has the same reduction by injectivity. Pair with $[S]e_0$ and use the inverse in \eqref{eq:display-cocycles}. It gives
\[
 [S]\log C_H=-C_p(Y,Y')+O((Y,Y')^{p^2}).
\]
Hence $p[SY_d^p]b=-1$, and exponentiation leaves this $S$-linear coefficient unchanged.

These same addition calculations imply $\log_G(S)=S+O(S^q)$ and
\[
 \log_H(Ye_d)=Ye_d+O(Y^{p^2}):
\]
compare the low-degree terms of the logarithm's additive identity on the indicated inputs. Therefore a rational bilinear change $\log_G(S)\sum_j a_j\log_{H,j}(\mathbf Y)$ contributes to neither of the two coefficients. An axes-normalized unit frame change adds or subtracts $\log I$. Its $Y$-linear and $S$-linear coefficients are integral, since higher powers of $I-1$ cannot contribute there. This proves the claimed invariance of their valuations.

For the full first homogeneous part, the diagonal display map $e_j\mapsto[\zeta^{p^j}]e_j$ commutes with the operators, including Frobenius on its coefficients. It sends $\sigma(S)$ to $\sigma(\zeta S)$, so gives the true $[\zeta]$ on the line. Its inverse dual sends $\tau(\mathbf Y)$ to $\tau(\zeta^{-p^j}Y_j)$. The evaluation pairing and the scalar-one quotient frame are invariant. As in DP37, averaging gives the unique invariant rational $b$.

Suppose an invariant mixed monomial has side degrees $a,b\ge1$ and total degree at most $p+1$. Then
\[
 a-\sum_jp^jb_j\equiv0\pmod{q-1},\qquad
 pb\le\sum_jp^jb_j\le p^db\le q.
\]
If the sum equals $a$, the degree bound forces $b=1$, $a=p$, and $j=1$. If the sum exceeds $a$, their difference can only be $q-1$; the upper bound forces $a=1$, $b=p$, and all dual factors to be $Y_d$. These are precisely the two monomials already computed, both of degree $p+1$. This proves \eqref{eq:primitive-first} and absence of smaller terms. Substitution of the inverse logarithms, whose linear parts are identities, preserves that first homogeneous part. The two stated factorial mixed derivatives follow.

For an integral function, the translation coefficient at source degree $p$ on the line is $D_G^p/p!$, since its logarithm has no nonlinear term of that degree. The corresponding last-dual coefficient is $D_d^p/p!$ by the last-axis calculation above. Both translation operators are integral, so $D_G^p$ and $D_d^p$ are divisible by $p$ on the entire integral function ring. This proves the asserted contrast with the two unit mixed jets. DP35 applied to a degree-$p+1$ coefficient of valuation $-1$ gives the numerical denominator rate.

Finally assume $h>2$. On the stated logarithmic path the first monomial in \eqref{eq:primitive-first} vanishes because $y_1=0$, leaving $-\epsilon^pt^{p+1}/p$. Put $D=D_G+\epsilon D_d$. Its two pure $p$th powers and every interior binomial term are divisible by $p$, so $D^p$ takes integral functions into $p\mathcal O[[S,\mathbf Y]]$. Consequently the integral $U$ factor contributes a multiple of $p$ to its $(p+1)$st derivative. There is no earlier positive term in the exponential gauge; its contribution at that order is
\[
 -U(0)\epsilon^p\frac{(p+1)!}{2p},
\]
a unit for odd $p$. It cannot cancel the first contribution. This proves \eqref{eq:primitive-path-unit}.
\end{proof}

At $h=2$ the path coefficient is $(\epsilon-\epsilon^p)/p$ and can vanish, which explains the restriction in \eqref{eq:primitive-path-unit}. For any prime, still with $h>2$, the same unit-jet proof works with $\exp(k\log B)$ for a unit scalar $k$ in place of the half-power. The asymptotic denominator rate also survives on a formal generic group-coordinate diagonal after an unramified extension with transcendental residue unit $\epsilon$: first restrict to $S,Y_1$, which retains the integral first coboundary and the coefficient $1/p$, then put $S=t,Y_1=\epsilon t$. In each nonzero homogeneous part the least-valuation residue is a nonzero polynomial in $\bar\epsilon$, so no cancellation of that minimum occurs. This proves the same rate there, but makes no claim for a prescribed algebraic CM parameter or for an arbitrary logarithmic diagonal.

